Иногда возникает логичное желание не только найти разделяющую прямую, но и постараться провести её на одинаковом
удалении от обоих классов, то есть максимизировать минимальный отступ

\subsection{Геометрическая интерпретация}

Это можно сделать, слегка поменяв функцию ошибки, а именно положив её равной:

\begin{gather*}
    F(M) = \max(0, 1 - M)\\
    L(w, x, y) = \lambda \| w \|_2^2 + \sum_i \max(0, 1 - y_i \langle w, x_i \rangle)\\
    \nabla_w L(w, x, y) = 2\lambda w + \sum_i
    \begin{cases}
        0, & 1 - y_i \langle w, x_i \rangle \leq 0 \\
        -y_i x_i, & 1 - y_i \langle w, x_i \rangle > 0
    \end{cases} \\
\end{gather*}
Почему же добавленная единичка приводит к желаемому результату?
Интуитивно это можно объяснить так: объекты, которые проклассифицированы правильно, но не очень \("\)уверенно\("\)
(то есть $0 \leq y_i (w,x_i) < 1$), продолжают вносить свой вклад в градиент и пытаются \("\)отодвинуть\("\) от себя
разделяющую плоскость как можно дальше.

\subsection{Оптимизационная задача для линейно разделимого случая}

К данному выводу можно прийти и чуть более строго; для этого надо совершенно по-другому взглянуть на выражение,
которое мы минимизируем.
Если мы максимизируем минимальный отступ, то надо максимизировать $\frac{2}{\abs{w}_2}$, то есть ширину
полосы при условии, что большинство объектов лежат с правильной стороны, что эквивалентно решению исходной задачи:
\[ \lambda \abs{w}_2^2 + \sum_i \max(0, 1 - y_i \langle w, x_i \rangle) \rightarrow min_w \]

Отметим, что первое слагаемое у нас обратно пропорционально ширине полосы, но мы и максимизацию заменили на
минимизацию, так что тут всё в порядке.
Второе слагаемое -- это штраф за то, что объекты неправильно расположены относительно разделительной полосы.
В конце концов, никто нам не обещал, что классы наши линейно разделимы и можно провести оптимальную плоскость вообще
без ошибок.

Итоговое положение плоскости задаётся всего несколькими обучающими примерами.
Это ближайшие к плоскости правильно классифицированные объекты, которые называют опорными векторами.
Весь метод, соответственно, зовётся методом опорных векторов, или support vector machine, или сокращённо SVM\@.
Начиная с шестидесятых годов это был сильнейший из известных методов машинного обучения.
В девяностые его сменили методы, основанные на деревьях решений, которые, в свою очередь, недавно передали «пальму
первенства» нейросетям.

Почему же SVM был столь популярен?
Из-за небольшого количества параметров и доказуемой оптимальности.
Сейчас для нас нормально выбирать специальный алгоритм под задачу и подбирать оптимальные гиперпараметры для этого
алгоритма перебором, а когда-то трава была зеленее, а компьютеры медленнее, и такой роскоши у людей не было.
Поэтому им нужны были модели, которые гарантированно неплохо работали бы в любой ситуации.
Такой моделью и был SVM\@.

Другие замечательные свойства SVM: существование уникального решения и доказуемо минимальная склонность к
переобучению среди всех популярных классов линейных классификаторов.
Кроме того, несложная модификация алгоритма, ядровый SVM, позволяет проводить нелинейные разделяющие поверхности.

\subsubsection{Двойственная задача}

\subsection{Неразделимый случай, slack-переменные}

В реальных задачах данные часто не являются линейно разделимыми из-за:
\begin{enumerate}
    \item Наложения классов (перекрывающиеся распределения)
    \item Выбросов (outliers) в данных
    \item Шума в измерениях
    \item Нелинейной границы между классами
\end{enumerate}
Подробнее на Рисунке~\ref{fig:unsep}
\begin{figure}[h]
    \centering
    \begin{tikzpicture}[scale=1.0]
        % Неразделимые данные
        \begin{scope}[xshift=0cm]
            \draw[->, thick] (-1,0) -- (3,0) node[right] {$x_1$};
            \draw[->, thick] (0,-1) -- (0,3) node[above] {$x_2$};
            
            % Точки класса +1
            \foreach \x/\y in {0.5/2, 1/2.2, 0.8/1.8, 1.5/2.5, 2/2.3, 2.5/2.7, 1/1.5, 1.2/1.3, 1.8/1.9}
            \fill[blue] (\x,\y) circle (2pt);
            
            % Точки класса -1
            \foreach \x/\y in {1.5/1.1, 1.8/0.8, 2/1.2, 2.5/0.7, 1.5/0.5, 2/0.3, 0.8/0.5, 1/0.8, 0.5/0.3}
            \fill[orange] (\x,\y) circle (2pt);
            
            % Выбросы
            \fill[blue] (0.2,0.3) circle (2pt) node[below] {выброс};
            \fill[orange] (2.8,2.5) circle (2pt) node[above] {выброс};
            
            % Попытка линейного разделения
            \draw[red, thick, dashed] (0,2.5) -- (2.5,0);
            
            \node at (1.5, -0.5) {Нелинейно разделимые данные};
        \end{scope}
    \end{tikzpicture}
    \caption{Пример неразделимых данных с выбросами}
    \label{fig:unsep}
\end{figure}

\underline{Slack-переменные (ослабляющие переменные)}

Для обработки неразделимого случая вводятся \textbf{slack-переменные} $\xi_i \geq 0$, которые позволяют объектам
нарушать запас на Рисунке~\ref{fig:slack}

\begin{align*}
    y_i(\bm{w}^\top \bm{x}_i + b) &\geq 1 - \xi_i, \quad i = 1, \dots, n \\
    \xi_i &\geq 0, \quad i = 1, \dots, n
\end{align*}

\begin{figure}[h]
    \centering
    \begin{tikzpicture}[scale=1.0]
        % Оси
        \draw[->, thick] (-1,0) -- (3,0) node[right] {$x_1$};
        \draw[->, thick] (0,-1) -- (0,3) node[above] {$x_2$};
        
        % Разделяющая гиперплоскость
        \draw[thick, red] (0,2) -- (2.5,0) node[midway, above left] {$\bm{w}^\top\bm{x} + b = 0$};
        
        % Границы запаса
        \draw[red, dashed] (0.5,2) -- (3,0.5) node[midway, above right] {$\bm{w}^\top\bm{x} + b = 1$};
        \draw[red, dashed] (-0.5,2) -- (2,-0.5) node[midway, below left] {$\bm{w}^\top\bm{x} + b = -1$};
        
        % Точки класса +1
        \foreach \x/\y in {1.5/1.8, 2/2, 1.2/1.5, 2.2/1.7}
        \fill[blue] (\x,\y) circle (2pt);
        
        % Точки класса -1
        \foreach \x/\y in {1.5/0.8, 1.8/0.5, 1.2/0.7, 1/1}
        \fill[orange] (\x,\y) circle (2pt);
        
        % Точки, нарушающие запас
        \fill[blue] (0.8,0.8) circle (2pt) node[below left] {$\xi_i > 1$};
        \fill[orange] (2,1.2) circle (2pt) node[above right] {$\xi_i > 0$};
        
        % Расстояния до правильной границы
        \draw[dashed, blue, ->] (0.8,0.8) -- (1.2,1.2);
        \draw[dashed, orange, ->] (2,1.2) -- (1.7,0.9);
        
        % Область правильной классификации
        \node[blue] at (2.2,2.2) {Класс +1};
        \node[orange] at (0.8,0.3) {Класс -1};
        
        % Подпись slack переменных
        \node at (1.5, -0.7) {$\xi_i$ измеряет степень нарушения запаса};
    \end{tikzpicture}
    \caption{Slack-переменные $\xi_i$ позволяют точкам нарушать запас}
    \label{fig:slack}
\end{figure}

\underline{Интерпретация slack-переменных}

\begin{itemize}
    \item $\xi_i = 0$: Точка $\bm{x}_i$ корректно классифицирована с запасом $\geq 1$.
    \item $0 < \xi_i \leq 1$: Точка $\bm{x}_i$ находится внутри запаса, но классифицирована правильно.
    \item $\xi_i > 1$: Точка $\bm{x}_i$ классифицирована неправильно
\end{itemize}

Величина $\xi_i$ измеряет, насколько точка $\bm{x}_i$ нарушает условие правильной классификации с запасом.

\underline{Задача оптимизации с slack-переменными}

Исходная задача SVM модифицируется

\begin{align*}
    \min_{\bm{w}, b, \bm{\xi}} &\quad \frac{1}{2} \|\bm{w}\|^2 + C \sum_{i=1}^n \xi_i \\
    \text{при условиях} &\quad y_i(\bm{w}^\top \bm{x}_i + b) \geq 1 - \xi_i, \quad i = 1, \dots, n \\
    &\quad \xi_i \geq 0, \quad i = 1, \dots, n
\end{align*}

где $C > 0$ --- гиперпараметр, задающий компромисс между:
\begin{itemize}
    \item Максимизацией запаса (минимизацией $\|\bm{w}\|^2$)
    \item Минимизацией ошибок (минимизацией $\sum \xi_i$)
\end{itemize}

\underline{Влияние параметра $C$}

\begin{itemize}
    \item \textbf{Большое $C$}: Большой штраф за ошибки → узкий запас, возможно переобучение
    \item \textbf{Малое $C$}: Маленький штраф за ошибки → широкий запас, возможна недотяжка
\end{itemize}
Проиллюстрированно на Рисунке~\ref{fig:C}
\begin{figure}[h]
    \centering
    \begin{tikzpicture}[scale=0.8]
        % Малое C
        \begin{scope}[xshift=0cm]
            \draw[->] (-1,0) -- (3,0);
            \draw[->] (0,-1) -- (0,3);
            
            % Данные
            \foreach \x/\y in {0.5/2, 1/2, 0.8/1.8, 1.5/2.2}
            \fill[blue] (\x,\y) circle (2pt);
            \foreach \x/\y in {1.5/1, 2/0.8, 1.8/0.5, 1/1.2}
            \fill[orange] (\x,\y) circle (2pt);
            
            % Большой запас (малое C)
            \draw[thick, red] (0,2.5) -- (3,0.8);
            \draw[red, dashed] (-0.5,2.8) -- (3.5,0.3);
            \draw[red, dashed] (0.5,2.2) -- (2.5,1.3);
            
            \node at (1.5, -0.5) {Малое $C$};
            \node at (1.5, -0.9) {Большой запас};
        \end{scope}
    \end{tikzpicture}
    \caption{Влияние параметра $C$ на ширину запаса}
    \label{fig:C}
\end{figure}

\subsection{Hinge loss}

\begin{definition}[Hinge loss]
    Функция потерь hinge (hinge loss) определяется как: \[ \mathcal{L}_{\text{hinge}}(y, z) = \max(0, 1 - yz) \]
    где $y \in \{-1, +1\}$ --- истинная метка, $z = \bm{w}^\top \bm{x} + b$ --- выход классификатора.
\end{definition}

\underline{Связь с задачей SVM}

Задачу SVM с slack-переменными можно переписать в терминах hinge loss:

\begin{align*}
    \min_{\bm{w}, b} &\quad \frac{1}{2} \|\bm{w}\|^2 + C \sum_{i=1}^n \max(0, 1 - y_i(\bm{w}^\top \bm{x}_i + b)) \\
    &= \min_{\bm{w}, b} \quad \frac{1}{2} \|\bm{w}\|^2 + C \sum_{i=1}^n \mathcal{L}_{\text{hinge}}(y_i, \bm{w}^\top
    \bm{x}_i + b)
\end{align*}

Таким образом, $\xi_i = \max(0, 1 - y_i(\bm{w}^\top \bm{x}_i + b))$.

\underline{Свойства hinge loss} изображены на Рисунке~\ref{fig:losses}
\begin{figure}[h]
    \centering
    \begin{tikzpicture}[scale=1.2]
        % Оси
        \draw[->, thick] (-2.5,0) -- (2.5,0) node[right] {$z = y \cdot (\bm{w}^\top\bm{x} + b)$};
        \draw[->, thick] (0,-0.5) -- (0,3.5) node[above] {Потери};
        
        % Hinge loss
        \draw[thick, blue] (-2.5,3.5) -- (1,0) -- (2.5,0);
        \node[blue, right] at (1.5,0.5) {Hinge loss: $\max(0, 1-z)$};
        
        % 0-1 loss
        \draw[thick, red] (-2.5,1) -- (0,1);
        \draw[thick, red] (0,0) -- (2.5,0);
        \draw[fill=white] (0,1) circle (2pt);
        \draw[fill=red] (0,0) circle (2pt);
        \node[red, right] at (2,0.5) {0-1 loss};
        
        % Log loss
        \draw[thick, green, dashed] plot[smooth, domain=-2.5:1.5] (\x, {ln(1+exp(-\x))/ln(2)});
        \node[green, right] at (2,2) {Log loss};
        
        % Вертикальная линия в z=1
        \draw[dashed] (1,0) -- (1,3);
        \node[below] at (1,-0.1) {$1$};
        \node[above] at (1,3) {Запас достигнут};
        
        % Области
        \node[align=center] at (-1.5, 0.5) {Неправильная\\классификация};
        \node[align=center] at (1.8, 0.5) {Правильная\\классификация};
    \end{tikzpicture}
    \caption{Сравнение hinge loss с другими функциями потерь}
    \label{fig:losses}
\end{figure}

\underline{Важные свойства hinge loss}

\begin{enumerate}
    \item \textbf{Выпуклость:} hinge loss выпукла, но не везде дифференцируема
    \item \textbf{Субградиент:} в точке $z = 1$ градиент плавает:
    \[ \partial \mathcal{L}_{\text{hinge}}(z) =
    \begin{cases}
        -1, & \text{если } z < 1 \\
        [ -1, 0 ], & \text{если } z = 1 \\
        0, & \text{если } z > 1
    \end{cases} \]
    
    \item \textbf{Разреженность:} при $z > 1$ градиент равен нулю → только точки на границе или внутри запаса
    влияют на решение (опорные векторы)
    \item \textbf{Робастность:} менее чувствительна к выбросам, чем log loss
\end{enumerate}

\underline{Оптимизация с hinge loss}

Задача SVM является задачей выпуклой оптимизации с ограничениями.

\underline{Прямая задача:}
\begin{gather*}
    \min_{\bm{w}, b, \bm{\xi}} \frac{1}{2} \|\bm{w}\|^2 + C \sum_{i=1}^n \xi_i\\
    \text{s.t. } y_i(\bm{w}^\top \bm{x}_i + b) \geq 1 - \xi_i, \quad \xi_i \geq 0\\
\end{gather*}

Решение выражается через \textbf{опорные векторы} (объекты с $\alpha_i > 0$):
\[ \bm{w} = \sum_{i=1}^n \alpha_i y_i \bm{x}_i \]

\subsection{Ядровой трюк, свойства ядер, примеры ядер}
Трюк изображён на Рисунке~\ref{fig:kernel}
\begin{figure}[h]
    \centering
    \begin{tikzpicture}[scale=1.0]
        % Нелинейно разделимые данные
        \begin{scope}[xshift=0cm]
            \draw[->] (-1.5,0) -- (1.5,0);
            \draw[->] (0,-1.5) -- (0,1.5);
            
            % Круговая граница
            \draw[gray, dashed] (0,0) circle (1);
            
            % Точки внутри круга (класс +1)
            \foreach \r/\a in {0.2/30, 0.3/150, 0.4/270, 0.5/60, 0.6/120, 0.7/240}
            \fill[blue] ({\r*cos(\a)}, {\r*sin(\a)}) circle (2pt);
            
            % Точки вне круга (класс -1)
            \foreach \r/\a in {1.2/20, 1.1/100, 1.3/200, 1.4/300, 1.1/50, 1.2/150, 1.3/250, 1.4/350}
            \fill[orange] ({\r*cos(\a)}, {\r*sin(\a)}) circle (2pt);
            
            \node at (0, -1.8) {Исходное пространство};
            \node at (0, -2.1) {Нелинейная граница};
        \end{scope}
        
        % Стрелка преобразования
        \begin{scope}[xshift=3cm]
            \node at (0,0) {$\xrightarrow{\phi(\bm{x})}$};
        \end{scope}
        
        % После преобразования
        \begin{scope}[xshift=6cm]
            \draw[->] (-1.5,0) -- (1.5,0) node[right] {$x_1^2$};
            \draw[->] (0,-1.5) -- (0,1.5) node[above] {$x_2^2$};
            
            % Линейная граница
            \draw[red] (0.5,-1.5) -- (0.5,1.5);
            
            % Точки после преобразования
            \foreach \r/\a in {0.2/30, 0.3/150, 0.4/270, 0.5/60, 0.6/120, 0.7/240}
            \fill[blue] ({0.5*(\r*cos(\a))^2}, {(\r*sin(\a))^2}) circle (2pt);
            
            \foreach \r/\a in {1.2/20, 1.1/100, 1.3/200, 1.4/300, 1.1/50, 1.2/150, 1.3/250, 1.4/350}
            \fill[orange] ({0.5*(\r*cos(\a))^2}, {(\r*sin(\a))^2}) circle (2pt);
            
            \node at (0, -1.8) {Признаковое пространство};
            \node at (0, -2.1) {Линейная граница};
        \end{scope}
    \end{tikzpicture}
    \caption{Идея ядрового метода: нелинейное преобразование в пространство большей размерности}
    \label{fig:kernel}
\end{figure}

\underline{Основная идея}

Вместо явного вычисления преобразования $\phi: \mathbb{R}^d \to \mathbb{R}^D$ (где $D \gg d$ или $D = \infty$),
мы работаем только со \textbf{скалярными произведениями}:
\[ K(\bm{x}, \bm{z}) = \langle \phi(\bm{x}), \phi(\bm{z}) \rangle \]

\underline{Определение ядра}

\begin{definition}[Положительно определённое ядро]
    Функция $K: \mathcal{X} \times \mathcal{X} \to \mathbb{R}$ называется \textbf{положительно определённым ядром},
    если:
    \begin{enumerate}
        \item $K$ симметрична: $K(\bm{x}, \bm{z}) = K(\bm{z}, \bm{x})$
        \item Для любого набора точек $\bm{x}_1, \dots, \bm{x}_n \in \mathcal{X}$ матрица Грама $[K(\bm{x}_i, \bm{x}_j)]_{i,j=1}^n$ положительно полуопределена
    \end{enumerate}
\end{definition}

\underline{Примеры ядер} изображены на Рисунке~\ref{tab:exmpls}

\begin{table}[h]
    \centering
    \begin{tabular}{|p{0.25\textwidth}|p{0.35\textwidth}|p{0.3\textwidth}|} \hline
        \textbf{Название ядра} & \textbf{Формула}                                           & \textbf{Параметры}                       \\ \hline
        Линейное               & $K(\bm{x}, \bm{z}) = \bm{x}^\top \bm{z}$                   & Нет                                      \\ \hline
        Полиномиальное         & $K(\bm{x}, \bm{z}) = (\gamma \bm{x}^\top \bm{z} + r)^d$    & $d \in \mathbb{N}, \gamma > 0, r \geq 0$ \\ \hline
        RBF (Гауссово)         & $K(\bm{x}, \bm{z}) = \exp(-\gamma \|\bm{x} - \bm{z}\|^2)$  & $\gamma > 0$                             \\ \hline
        Сигмоидное             & $K(\bm{x}, \bm{z}) = \tanh(\gamma \bm{x}^\top \bm{z} + r)$ & $\gamma > 0, r$                          \\ \hline
    \end{tabular}
    \caption{Часто используемые ядерные функции}
    \label{tab:exmpls}
\end{table}

\underline{Математические свойства ядер}

\begin{theorem}[Теорема Мерсера]
    Если $K$ --- непрерывное симметричное положительно определённое ядро на компактном множестве $\mathcal{X}
    \subset \mathbb{R}^d$, то существует последовательность $\{\lambda_i\}_{i=1}^\infty$ ($\lambda_i \geq 0$) и
    ортонормированная система $\{\phi_i\}_{i=1}^\infty$ в $L^2(\mathcal{X})$ такая, что:
    \[ K(\bm{x}, \bm{z}) = \sum_{i=1}^\infty \lambda_i \phi_i(\bm{x}) \phi_i(\bm{z}) \]
    Такой ряд сходится абсолютно и равномерно.
\end{theorem}

\underline{Замкнутость относительно операций}

\begin{enumerate}
    \item \textbf{Сумма ядер:} Если $K_1$ и $K_2$ --- ядра, то $K(\bm{x}, \bm{z}) = K_1(\bm{x}, \bm{z}) +
    K_2(\bm{x}, \bm{z})$ --- ядро
    \item \textbf{Произведение ядер:} Если $K_1$ и $K_2$ --- ядра, то $K(\bm{x}, \bm{z}) = K_1(\bm{x}, \bm{z})
    \cdot K_2(\bm{x}, \bm{z})$ --- ядро
    \item \textbf{Масштабирование:} Если $K$ --- ядро и $c > 0$, то $cK$ --- ядро
    \item \textbf{Композиция с функцией:} Если $K$ --- ядро и $f: \mathbb{R} \to \mathbb{R}$ --- многочлен с
    неотрицательными коэффициентами, то $f(K)$ --- ядро
\end{enumerate}

\underline{Настройка гиперпараметров:}

Для SVM с RBF ядром нужно настраивать:
\begin{itemize}
    \item $C$: компромисс между запасом и ошибками.
    \item $\gamma$: сложность границы решения
\end{itemize}